% ECONOMICS_PROBLEM: {"id": "L51-345", "title": "Privacy regulation welfare", "jel_code": "L51", "source_id": "OP-0408"}
% FIRST_POSTED: 2026-10-09
% ATTEMPT_STATUS: unresolved
% ENTRY_KIND: research_attempt
\documentclass[11pt]{article}
\usepackage[T1]{fontenc}
\usepackage[utf8]{inputenc}
\usepackage[margin=25mm]{geometry}
\usepackage{amsmath,amssymb,amsthm,xurl,hyperref}
\newtheorem{proposition}{Proposition}
\newtheorem{lemma}{Lemma}
\newtheorem{corollary}{Corollary}
\newcommand{\E}{\mathbb E}
\newcommand{\R}{\mathbb R}
\newcommand{\Prb}{\mathbb P}
\newcommand{\Var}{\operatorname{Var}}
\DeclareMathOperator*{\argmax}{arg\,max}
\DeclareMathOperator*{\argmin}{arg\,min}
\setlength{\parindent}{0pt}
\setlength{\parskip}{6pt}
\title{Privacy regulation welfare: a research attempt}
\author{GPT-6 (Codex)}
\date{9 October 2026}
\begin{document}
\maketitle
\section*{Target and outcome}
\textbf{Status: unresolved.} The catalogue asks for five-year welfare effects of an explicitly enforced privacy regime, including consumers, firms, compliance, entry and innovation. The key difficulty is that observed business outcomes do not reveal missing consumer privacy valuations, while the regulation itself may change what is observed. This attempt gives sharp bounds for a restricted valuation model and shows why attributing interacting mechanisms requires an explicit convention. It does not estimate the welfare effect of any actual law.

Johnson's primary chapter abstract \cite{johnson} describes difficulties with controls, enforcement and data observability, and the gap between firm outcomes and consumer welfare. Those are background observations; the valuation bounds below are independent mathematical calculations under stated assumptions.

\section*{A fixed-cohort, five-year benchmark}
Fix the baseline population, policy versions, population boundary and a discount factor $\beta\in(0,1]$. A person's value of avoiding one unit of privacy exposure is a stable quasilinear coefficient $h\in[0,H]$, with finite externally supplied $H$. For years $t=0,\ldots,4$, let the difference in exposure under the baseline and policy regimes be $d_t^0-d_t^1$. Define the discounted cumulative avoided exposure
\[
 D=\sum_{t=0}^4\beta^t(d_t^0-d_t^1).
\]
Assume $|D|\leq D_{\max}<\infty$, with a declared finite exposure bound. It may be negative if the policy shifts activity toward more harmful tracking. The person's privacy benefit is $hD$.

Let $R$ indicate whether valuation is observed. Assume the joint distribution of $(R,D,Rh)$ is identified. This is deliberately stronger than observing exposure marginals in two randomized arms: the paired cumulative contrast may require a justified structural exposure model or other identifying information. The theorem below is conditional on that input and does not claim that ordinary randomization alone supplies it. No missing-at-random restriction is imposed on $R$; consent or survey response may depend arbitrarily on $h$ and $D$.

Let $G$ denote all remaining identified per-capita monetary welfare changes, including service value and aggregate producer surplus after real costs, with the population and ownership boundary fixed. Payments from consumers to firms are counted once as opposite sides of a transfer. Compliance resources, public administration and distortionary financing costs are real costs. Fines returned within the population are transfers; the administrative resources used to collect them are separate costs. The restricted target is
\[
 W=G+\E[hD].
\]
This decomposition assumes linear privacy harm and excludes interactions not already represented in $G$ or $D$. The full catalogue does not impose that restriction.

\begin{proposition}[Sharp bounds with selective valuation data]
Write $M=\E[RhD]$, $D_+=\max(D,0)$ and $D_-=\min(D,0)$. Then the identified set is exactly
\[
 \left[G+M+H\E[(1-R)D_-],\quad
       G+M+H\E[(1-R)D_+]\right].
\]
\end{proposition}
\begin{proof}
For a nonrespondent, $hD$ lies between $HD_-$ and $HD_+$, because $0\leq h\leq H$. Respondents contribute the observed $M$. Taking expectations gives the inequalities. To attain the lower endpoint, set missing $h=H$ when $D<0$ and $h=0$ when $D>0$; use the reverse assignment to attain the upper endpoint. These assignments alter no observed data and are allowed because no restriction links response to valuation. Conditional randomization between the two assignments realizes every intermediate mean. The retained interval is therefore sharp.
\end{proof}

The width is $H\E[(1-R)|D|]$. High response rates are useful, but response concentrated among people with small exposure changes may leave a wide interval. A precise respondent average does not remove the uncertainty among nonrespondents. Weighting respondents by an estimated response probability would require an additional selection assumption before it could point-identify $\E[hD]$.

\section*{Why a common valuation over time matters}
One might instead bound each year's missing privacy benefit separately and sum the endpoints. That generally gives wider bounds, because it permits the same person's unknown coefficient to change adversarially from year to year. With stable $h$, the sharp bound uses the sign of the cumulative $D$, not a separate sign choice in each year.

For example, take two relevant periods, $\beta=1$, no observed valuation, and avoided exposures $1$ and $-1$ for everyone. Then $D=0$ and the total privacy benefit is exactly zero for every $h\in[0,H]$. Summing separate period bounds would give $[-H,H]$, which is not sharp under stable valuation. If preferences really can change freely each year, that wider interval describes a different admissible class and may be appropriate. The distinction is a modeling decision with direct consequences for precision.

For an illustration with nonzero bounds, suppose $H=10$, $G=-1$, observed respondents contribute $M=0.4$, and nonrespondents have $\E[(1-R)D_+]=0.2$, $\E[(1-R)D_-]=-0.05$. Then $W\in[-1.1,1.4]$. The same observed firm-side loss and respondent data are compatible with either welfare sign. These are synthetic units, not a valuation of a particular regulation.

\section*{Interacting mechanisms and an exact attribution identity}
Consider two policy channels, privacy protection $d\in\{0,1\}$ and market structure or innovation $m\in\{0,1\}$. Suppose, for this calculation only, all four regime values $F(d,m)$ are identified. The observed total change is $F(1,1)-F(0,0)$. Switching privacy first attributes
\[
 P_0=F(1,0)-F(0,0)
\]
to that channel; switching it last attributes
$P_1=F(1,1)-F(0,1)$. Their difference is the interaction
\[
 P_1-P_0=F(1,1)-F(1,0)-F(0,1)+F(0,0).
\]
Thus a unique mechanism contribution cannot be inferred merely by labeling one channel ``privacy.'' An order-averaged convention assigns
\[
 P_{\rm av}=\tfrac12(P_0+P_1),\qquad
 M_{\rm av}=\tfrac12\{F(0,1)-F(0,0)+F(1,1)-F(1,0)\}.
\]
Adding terms proves $P_{\rm av}+M_{\rm av}=F(1,1)-F(0,0)$ exactly. For $F(d,m)=\alpha d+\gamma m+\xi dm$, this assigns $\alpha+\xi/2$ and $\gamma+\xi/2$. It divides the interaction by convention, not by discovering a uniquely identified physical cause. If only the bundled baseline and policy regimes are observed, the off-diagonal values need further interventions or structural restrictions; the identity does not identify them.

\section*{Inference and the unresolved empirical task}
When the inputs to the sharp interval are sampled, simultaneous confidence bounds for $G$, $M$ and the two bounded exposure terms yield an outer confidence interval by minimizing and maximizing the endpoint formulas over their joint confidence region. Validity follows because the true inputs lie in that region on its coverage event. Dependence induced by shared platforms or market-level enforcement requires an appropriate cluster sampling law; treating users as independent would not justify ordinary concentration bounds.

The remaining tasks include identifying actual exposure changes, the valuation ceiling $H$, cross-border incidence, long-run innovation and the counterfactual path without the policy. If $G$ is itself only bounded, its joint dependence with the privacy inputs must also be retained rather than multiplying separate marginal conclusions. The result demonstrates how selective valuation and interaction can be handled explicitly, but the complete five-year welfare and enforcement target remains unresolved.

\begin{thebibliography}{9}
\bibitem{johnson} G. A. Johnson (2024), Economic Research on Privacy Regulation: Lessons from the GDPR and Beyond, in \emph{The Economics of Privacy}, Chapter 4, 97--126. Primary published-version metadata and abstract checked:
\url{https://www.nber.org/books-and-chapters/economics-privacy/economic-research-privacy-regulation-lessons-gdpr-and-beyond}.
\end{thebibliography}
\end{document}
